\clearpage
\renewcommand{\sectioncolor}{sand}
\section{``The next Euler-like synthesis'' --- assessed}\label{sec:euler}
\markboth{The next synthesis}{}

The pasted text closes on an attractive idea: that just as Euler united arithmetic, geometry,
trigonometry, calculus and complex numbers, a future synthesis might unite
\emph{matter, information, probability, geometry and life}. It is worth taking seriously enough to
test --- and the test comes from Part VI itself.

\begin{center}
\begin{tikzpicture}[font=\footnotesize,
  bx/.style={rounded corners=4pt,draw=#1,line width=1.2pt,fill=#1!7,inner sep=9pt,
             text width=7.45cm,align=left,minimum height=5.6cm}]
 \node[bx=bio] (a) at (-4.35,0)
   {\textbf{\color{bio}\normalsize What made Euler's synthesis real}\\[5pt]
    It was \textbf{not} a list of fields joined by ``$+$''.\\[4pt]
    It was a specific claim: two conceptually different \emph{rules} determine the same set of
    ordered pairs, because both\\[3pt]
    \quad$\bullet$ map sums to products\\
    \quad$\bullet$ change in proportion to size\\
    \quad$\bullet$ are periodic\\
    \quad$\bullet$ are self-regulating\\[4pt]
    {\scriptsize Four checkable properties. A synthesis is a \textbf{shared entailment}, not a shared
    topic.}};
 \node[bx=crimson] (b) at (4.35,0)
   {\textbf{\color{crimson}\normalsize What the proposed synthesis currently is}\\[5pt]
    $\boxed{\text{Matter}+\text{Information}+\text{Probability}+\text{Geometry}+\text{Life}}$\\[6pt]
    A list of \emph{subject areas}, with no shared entailment named, no rules compared, and no
    property that both sides of an equation would have to possess.\\[5pt]
    {\scriptsize This is a research \textbf{aspiration}, and a good one. It is not yet a synthesis,
    and calling it one by drawing a box around it is precisely the move Part VI teaches you to
    distrust.}};
\end{tikzpicture}
\end{center}
\vspace{-2pt}
\begin{center}\begin{minipage}{0.93\textwidth}\footnotesize\color{slate}
\textbf{Figure 4.}\; The test any claimed synthesis must pass. Euler's passes it; the proposed one
does not yet, and saying so is not a criticism of the ambition --- it is a specification of what
would have to be produced.
\end{minipage}\end{center}

\subsection{What a real candidate would look like}

\begin{movebox}[title={The shape of a genuine synthesis in this domain}]
To count, a proposed unification would have to exhibit \textbf{one structure with two conceptually
different descriptions that provably entail the same properties}. Candidates that already have this
shape, at small scale:
\begin{center}\small
\begin{tabular}{@{}>{\raggedright\arraybackslash}p{5.2cm}>{\raggedright\arraybackslash}p{5.2cm}>{\raggedright\arraybackslash}p{5.5cm}@{}}
\toprule
\rowcolor{bio!13}
\textbf{\color{bio}Description A} & \textbf{\color{bio}Description B} & \textbf{\color{bio}Shared entailment}\\
\midrule
reaction network as a \emph{graph} & reaction network as a \emph{dynamical system} &
 deficiency theorems: graph topology $\Rightarrow$ existence, uniqueness, stability of equilibria\\
irreversible pulling work & equilibrium free-energy difference &
 Jarzynski: $\langle e^{-\beta W}\rangle = e^{-\beta\Delta F}$\\
accuracy of a molecular copier & energy dissipated per step &
 kinetic proofreading: fidelity is bought with dissipation\\
\rowcolor{bio!6}
precision of a biochemical clock & entropy production rate &
 thermodynamic uncertainty relation\\
\bottomrule
\end{tabular}
\end{center}
\vspace{2pt}
Each of these \emph{is} a small Euler: a quantity computed two incompatible ways, provably equal.
\textbf{The realistic path to a large synthesis is to accumulate more of these, not to name the
fields you hope will merge.}
\end{movebox}
